\begin{figure}[htbp]
\centering
\begin{tikzpicture}[font=\small,>=Stealth]

\filldraw[fill=blue!8,draw=blue] (0,0) rectangle (1.8,1.3);
\draw[->] (0,-.2)--(1.8,-.2) node[midway,below] {$\Delta u$};
\draw[->] (-.2,0)--(-.2,1.3) node[midway,left] {$\Delta v$};
\node at (.9,2) {parameter rectangle};
\draw[->,thick] (2.1,.7)--(3,.7) node[midway,above] {$\sigma$};
\filldraw[fill=blue!8,draw=blue,thick] (3.3,0) .. controls (4,.4) and (4.7,.2) .. (5.3,.5)
 .. controls (5.2,1.1) and (5.6,1.4) .. (5.3,1.7)
 .. controls (4.7,1.4) and (4,1.8) .. (3.3,1.3)
 .. controls (3.5,.8) and (3.2,.5) .. (3.3,0);
\node at (4.3,2) {curved patch};
\draw[->,thick] (5.7,.7)--(6.6,.7);
\filldraw[fill=blue!8,draw=blue] (7,0)--(9,0)--(9.6,1.3)--(7.6,1.3)--cycle;
\draw[->,very thick] (7,0)--(9,0) node[midway,below] {$\sigma_u\Delta u$};
\draw[->,very thick] (7,0)--(7.6,1.3) node[pos=.65,right] {$\sigma_v\Delta v$};
\node at (8.3,2) {tangent parallelogram};
\node at (5,-1.25) {area to first order: $\|\sigma_u\times\sigma_v\|\Delta u\Delta v$};

\end{tikzpicture}
\caption{The derivative replaces a small curved patch by its tangent parallelogram. The cross-product norm gives the local scalar area factor.}
\label{fig:surface-area-mechanism}
\end{figure}
